
% BEGIN bib
 
\par\medskip\noindent\textbf{References}\par
  
\par\noindent  [1]  Dimitri Ara, Albert Burroni, Yves Guiraud, Philippe Malbos, François Métayer, and Samuel Mimram.  Polygraphs: from rewriting to higher categories.  Cambridge University Press, 2025.  
 
\par\noindent  [2]  Franz Baader and Tobias Nipkow.  Term rewriting and all that.  Cambridge University Press, 1998.  
 
\par\noindent  [3]  Christopher J Dean, Eric Finster, Ioannis Markakis, David Reutter, and Jamie Vicary.  Computads for weak $\omega$-categories as an inductive type.  Advances in Mathematics, 450, 2024.  
 
\par\noindent  [4]  Yves Guiraud and Philippe Malbos.  Identities among relations for higher-dimensional rewriting systems.  In Operads 2009, volume 26 of Séminaires et congrès, pages 145–161. Société Mathématique de France, 2013.  
 
\par\noindent  [5]  Yves Guiraud and Philippe Malbos.  Polygraphs of finite derivation type.  Mathematical Structures in Computer Science, 28(2):155–201, 2018.  doi:10.1017/S0960129516000220.  
 
\par\noindent  [6]  Allen Hatcher.  Algebraic Topology.  Cambridge University Press, 2002.  
 
\par\noindent  [7]  Simon Henry and Samuel Mimram.  Tietze equivalences as weak equivalences.  Applied Categorical Structures, 30(3):453–483, 2022.  doi:10.1007/S10485-021-09662-W.  
 
\par\noindent  [8]  Donald E. Knuth and Peter B. Bendix.  Simple word problems in universal algebras.  In Computational problems in abstract algebra, pages 263–297. Pergamon, 1970.  
 
\par\noindent  [9]  Nicolai Kraus and Jakob von Raumer.  Path spaces of higher inductive types in homotopy type theory.  In 2019 34th Annual ACM/IEEE Symposium on Logic in Computer Science (LICS), pages 1–13. IEEE, 2019.  doi:10.1109/LICS.2019.8785661.  
 
\par\noindent  [10]  Nicolai Kraus and Jakob von Raumer.  A rewriting coherence theorem with applications in homotopy type theory.  Mathematical Structures in Computer Science, 32(7):982–1014, 2022.  doi:10.1017/S0960129523000026.  
 
\par\noindent  [11]  Yves Lafont.  A new finiteness condition for monoids presented by complete rewriting systems (after Craig C. Squier).  Journal of Pure and Applied Algebra, 98(3):229–244, 1995.  
 
\par\noindent  [12]  Yves Lafont, François Métayer, and Krzysztof Worytkiewicz.  A folk model structure on omega-cat.  Advances in Mathematics, 224(3):1183–1231, 2010.  
 
\par\noindent  [13]  Samuel Mimram.  Polygraphs in Agda.  Software, swhId: swh:1:dir:27f2a0c3a496b72c3f4cd458ab51ea174f448eb7 (visited on 2025-06-23).  URL: https://github.com/smimram/agda-polygraphs, doi:10.4230/artifacts.23666.  
 
\par\noindent  [14]  Samuel Mimram.  Categorical coherence from term rewriting systems.  In Marco Gaboardi and Femke van Raamsdonk, editors, 8th International Conference on Formal Structures for Computation and Deduction (FSCD), volume 260 of LIPIcs, pages 16:1–16:17. Schloss Dagstuhl, 2023.  doi:10.4230/LIPICS.FSCD.2023.16.  
 
\par\noindent  [15]  Maxwell Herman Alexander Newman.  On theories with a combinatorial definition of “equivalence”.  Annals of mathematics, 43(2):223–243, 1942.  
 
\par\noindent  [16]  Kristina Sojakova, Floris van Doorn, and Egbert Rijke.  Sequential colimits in homotopy type theory.  In Holger Hermanns, Lijun Zhang, Naoki Kobayashi, and Dale Miller, editors, LICS ’20: 35th Annual ACM/IEEE Symposium on Logic in Computer Science, pages 845–858. ACM, 2020.  doi:10.1145/3373718.3394801.  
 
\par\noindent  [17]  Craig C Squier, Friedrich Otto, and Yuji Kobayashi.  A finiteness condition for rewriting systems.  Theoretical Computer Science, 131(2):271–294, 1994.  doi:10.1016/0304-3975(94)90175-9.  
 
\par\noindent  [18]  Heinrich Tietze.  Über die topologischen Invarianten mehrdimensionaler Mannigfaltigkeiten.  Monatshefte für Mathematik und Physik, 19:1–118, 1908.  
 
\par\noindent  [19]  The Univalent Foundations Program.  Homotopy Type Theory: Univalent Foundations of Mathematics.  https://homotopytypetheory.org/book, Institute for Advanced Study, 2013.  
 
\par\noindent  [20]  Andrea Vezzosi, Anders Mörtberg, and Andreas Abel.  Cubical Agda: a dependently typed programming language with univalence and higher inductive types.  Proceedings of the ACM on Programming Languages, 3:1–29, 2019.  doi:10.1145/3341691.  
  
% END bib
